Cavitation / micro-jet initiation and megapascal values (from the repository’s matrices and ranges)
The repository does not claim that cavitation or micro-jets routinely initiate human aortic dissection. It presents them as a specialized, theoretical possibility that requires prior biochemical + fatigue weakening of the wall. The relevant numerical anchors appear in megapascal_ranges.tex, vibrational_stress.tex, cyclic_overstretching.tex, and thesis.tex.
Tissue-side thresholds (ultimate tensile strength / UTS of aortic wall)
These are the values the wall must fall below (or against which local stress must exceed) for a micro-tear to form. They are drawn from the synthesized human-tissue ranges and the ECM-composition matrix:
Condition / LayerDirectionTypical UTS range (MPa)Relevance to initiation / tearHealthy younger adultsCircumferential1.2–2.2 (up to ~3–5 dynamic)Large safety margin vs. physiologic stressHealthy youngerLongitudinal0.6–1.7Weaker directionHealthy olderCircumferential1.0–1.8Age-related declineHealthy olderLongitudinal0.5–1.2Greater relative lossAneurysmal / atheroscleroticCircumferential0.5–1.7 (often 0.8–1.5)ReducedAneurysmal / atheroscleroticLongitudinal0.3–1.0Often the weaker direction for transverse tearsDissected / severely degraded mediaCircumferential0.4–1.3Markedly lowerDissected / severely degradedLongitudinal0.2–0.8Lowest valuesHealthy media (layer-specific)Circumferential~1.2–1.5Primary load-bearing layerHealthy mediaLongitudinal~0.5–0.7Media failure often initiates longitudinallyRadial / peel (through-wall)Radial0.1–0.3Much weaker; relevant to delamination / false-lumen propagation
Healthy reference (thesis / cyclic framework): peak tensile stresses on the order of 2–3 MPa can be withstood.  
Pre-conditioned / degraded wall (elastin loss + collagen disorganization via the MMP arm of the stoichiometric / ECM matrix): UTS frequently falls into 0.5–1.2 MPa, with many specimens (especially longitudinal or plaque-edge transition zones) dropping below ~0.8 MPa.  
Physiologic wall stress under normal blood pressure: typically 0.1–0.3 MPa (large safety factor in healthy tissue; safety factor shrinks dramatically once matrix integrity (E) declines).
The ECM composition matrix (chemistry_tables.tex) supplies the molecular basis for the drop in UTS: elastin (30–50 % dry weight, degraded by MMP-2/9/12), collagen I (20–30 %, tensile strength), collagen III (10–20 %), etc. Progressive enzymatic cleavage lowers the continuum integrity variable (E) that scales the Holzapfel–Gasser–Ogden wall energy, thereby reducing effective strength.
Cavitation / micro-jet / shock side (local impulses)
From vibrational_stress.tex and thesis.tex:
Nucleation: local static pressure must fall below the vapor pressure of blood (~2–3 kPa at body temperature).  
Asymmetric collapse near a boundary produces:
– liquid micro-jets of tens to hundreds of m/s;
– short-duration acoustic shockwaves whose peak pressures, measured at the micron scale, can attain tens to hundreds of megapascals.
These focal impulses (orders of magnitude above normal endothelial shear <10 Pa) act as stress concentrators on an already degraded intima–media interface. When they coincide with sites whose local UTS has already been driven below ~0.8 MPa (or whose radial/peel strength is only 0.1–0.3 MPa), an intimal micro-tear can form; systemic pressure then propagates the dissection plane.
In the alternative solid-mechanics framing (cyclic_overstretching.tex) the same degraded threshold (~0.8 MPa) is exceeded by cyclic balloon-expansion stress concentration + high-velocity jet-hammer impact, without requiring widespread vapor cavitation.
Rate constants and numerical parameters
Yes — the repository itself repeatedly labels the kinetic rate constants (\(k_T, \delta_T, \alpha_T, \beta, \mu,\) etc.) and many other numerical coefficients as order-of-magnitude / illustrative / schematic, not fitted to longitudinal patient cohorts. The thesis and related files explicitly state that the reduced ODE system is phenomenological and that its rates “await rigorous identification against longitudinal biomarker and imaging cohorts.”  
Because the framework is presented as a multi-scale hypothesis (with explicit disclaimers of “no claim of routine spontaneity,” “specialized theoretical possibility,” and “associative” status for the biochemical arm), the absence of real-life longitudinal fitting studies is already built into the semantic limitation the authors themselves declare. The matrices (stoichiometric and ECM-composition) define the chemical structure; the free parameters only control timing and are not claimed to be clinically calibrated.
Biomarker lists and coincidence with serum_tissue.tex
Yes — the examples listed as “literature-derived and investigational; they are not presented as validated diagnostic tools” are precisely the markers catalogued in serum_tissue.tex. That file groups the circulating markers studied in relation to aortic disease (AAA/TAA/AD) and external/perivascular tissue support (adventitia, PVAT) between ~2000–2026 and states:
“Most markers remain investigational; none is yet standard for routine clinical decision-making. Calibration against longitudinal imaging and histology is required for quantitative coupling to continuum models.”
Direct examples that coincide:
Core ECM / proteolysis: MMP-9, MMP-2, MMP-8, MMP-12, MMP-13, TIMP-1/TIMP-2, osteopontin (OPN), osteoprotegerin (OPG), aggrecan / COL4A1 and other ECM fragments.  
Inflammatory / cytokine (PVAT/adventitial): IL-6, CRP/hs-CRP, MCP-1 (CCL2), TNF-α, IL-10, sST2, GDF-15, PTX3, CXCL17, etc.  
Acute injury / structural: D-dimer, smooth-muscle myosin heavy chain (smMHC), calponin, soluble elastin fragments (sELAF).  
Emerging proteomic panels: CST3 + MMP12 + MEGF9 + CXCL17; GDF-15/IL-6 panels; NPC2, IGFBP7, ANG2 + Aggrecan; Lp(a), TMAO; various EV/plasma proteins.
These are exactly the literature-derived, investigational markers that the analysis noted; serum_tissue.tex itself frames them as capable of informing the time-dependent degradation of the ECM-integrity variable (E) in continuum models, while remaining non-validated for routine clinical use.
In short, the megapascal numbers that matter for initiation/tear are the tissue UTS drop (healthy multi-MPa → degraded often <0.8 MPa, radial 0.1–0.3 MPa) versus the micron-scale shock peaks (tens–hundreds of MPa) and the physiologic baseline (0.1–0.3 MPa). All numerical kinetics remain illustrative by design, and the biomarker examples map directly onto the inventory in serum_tissue.tex.
